\documentclass[10pt]{article}
\usepackage[margin=20mm]{geometry}
\usepackage{amsmath,amssymb,fontspec,unicode-math,xurl,hyperref}
\IfFontExistsTF{DejaVu Serif}{\setmainfont{DejaVu Serif}}{\setmainfont{Latin Modern Roman}}
\setmathfont{latinmodern-math.otf}
\setlength{\parindent}{0pt}\setlength{\parskip}{6pt}\sloppy
\title{Technical Verification: Neural Encoding and Decoding at Scale}
\author{AI-assisted technical review}
\date{}
\begin{document}
\maketitle
Abstract: The conditional-prediction objective is well motivated, but the pinned public release differs from the paper in within-modality input masking and in the interpretation of continuous-behavior R-squared. Its unimodal output modules are also omitted from normal optimizer and checkpoint traversal. Independent finite calculations and a CPU PyTorch toy validate these mechanisms and limited repairs. Historical benchmark effects remain unverified, and the paper’s empirical strengths and narrower supported interpretations are retained. Adjudication remains pending.\par
\section*{Paper identity and reviewed version}
This report concerns Neural Encoding and Decoding at Scale, by Yizi Zhang et al., ICML 2025, PMLR volume 267; identifier `I25-43679`, OpenReview `vOdz3zhSCj`. The reviewed object is the frozen published proceedings PDF, retrieved on 25 October 2026. Page references below are one-based PDF pages, which coincide with the internal page numbers. They are not the continuous PMLR volume pagination.\par
\section*{Scope and what the objective supports}
NEDS trains a shared transformer with neural, behavioral, within-modality and cross-modal reconstruction tasks, then adapts session-specific components to new animals. Main Sections 1–7 and Appendices A–H were read. Relevant architecture, masking, optimizer, checkpoint, data-split and metric paths were statically traced in the frozen public NEDS source at commit `61d2ef65556dbbc2b828d542f19bd4e7134c92f3`, dated 20 March 2026. This later release is not presumed identical to the code that generated the 2025 results. No author module, training program or model checkpoint was executed. Independent calculations use only small arrays and a newly written CPU PyTorch example.\par
The core statistical motivation is sound. If a Poisson prediction has rate \(r>0\) and target conditional mean \(\mu\), its expected negative log-likelihood, apart from a constant, is \[r-\mu\log r,\qquad \frac{\partial}{\partial r}(r-\mu\log r)=1-\mu/r.\] For \(\mu>0\) it is minimized at \(r=\mu\), even when the true count distribution is not Poisson. MSE similarly targets a conditional mean; categorical cross-entropy targets class probabilities. A finite check with equally likely counts \(0,0,6\) confirms the Poisson quasi-loss minimum at mean 2 among the tested rates. The code uses log-rate Poisson loss and exponentiates predictions for evaluation, resolving the positive-rate link only implicit in Equation (1).\par
These objectives justify conditional prediction under the actual input masks. They do not by themselves establish a mutually compatible joint generative distribution, causal influence of behavior, or transfer without fine-tuning. The animal-level pretraining holdout is a meaningful strength. The reported evaluation then adapts on 70\% of each held-out animal’s trials, selects settings on 10\% and tests on 20\%. The inspected main data-preparation path implements this random trial partition and excludes listed test sessions from multi-session pretraining. Exact identifiers and sample membership were not independently compared with raw data.\par
\section*{The released within-modality task uses additional information}
Problem 1 (The source implements a different conditional reconstruction task).  Section 3.2 and Figure 1 define within-modality reconstruction as \(\mathbb E[X_{\rm masked}\mid X_{\rm visible}]\) or its behavioral counterpart, with the other modality absent. The pinned implementation leaves the other modality visible.\par
In `src/trainer/base.py`, lines 100–106 set a random neural mask for `self-spike` and zero masks for all behavioral modalities; `self-behavior` analogously leaves spikes unmasked. Zero means visible, as confirmed by `src/multi\_modal/mm.py`, lines 152–160. All multimodal inputs are concatenated and attend to one another. Hence neural reconstruction can use \(\mathbb E[X_{\rm masked}\mid X_{\rm visible},Y]\). The mixed-batch path at lines 301–315 of the same model file has the same information exposure.\par
This changes the task mathematically. Let a masked neural value \(X\) be a uniform bit, let visible same-modality information be constant, and let behavior satisfy \(Y=X\). The intended within-modality optimal MSE is \(1/4\); with visible behavior it is zero. The independent two-outcome calculation verifies both values. Thus the two masks cannot be treated as equivalent regularizers or as identifying the same conditional expectation. This is a Major method-to-source discrepancy, particularly for interpreting the masking ablation in Table 2. It is not test-set leakage: it concerns information supplied during an explicitly multimodal training task.\par
The source also uniformly samples five named tasks, separating the two within-modality variants, rather than the four categories drawn uniformly in Equation (1). Their aggregate probability is consequently \(2/5\), not \(1/4\), while encoding and decoding each receive \(1/5\). To implement the paper literally, input visibility must be separated from the loss-selection mask: hide the other modality completely but compute reconstruction loss only for the selected within-modality targets. The source already distinguishes these concepts for selected-variable encoding. Alternatively, revise the stated conditionals and report the actual sampling weights. The finite example verifies the information-set distinction; neither a patched training run nor a historical mask reconstruction was performed.\par
\section*{Unimodal output heads are not registered model components}
Problem 2 (The pinned unimodal baseline omits its output modules from optimization and checkpoints).  The unimodal setup stores session-specific output decoders and temporal pooling weights inside ordinary Python dictionaries. They are not discovered by the model’s parameter and state traversal.\par
In `src/multi\_modal/mm.py`, lines 114–124 create ordinary dictionaries for the output decoders and temporal pooling weights, then insert decoder modules and nested parameter dictionaries. These components are used in the forward path at lines 253–277. But `src/train.py`, lines 297–302 constructs AdamW solely from `model.parameters()`, and `src/trainer/base.py`, lines 544–561 saves `model.state\_dict()`. No manual addition of those dictionary-held components was found along these paths. By contrast, multimodal encoders are in a registered module dictionary, and their output modules and pooling parameters are registered attributes.\par
An independent PyTorch model reproduces the container mechanism without importing any NEDS code. Its dictionary-held output head has gradient norm \(129\), yet a real AdamW step changes its weight by exactly zero. Its pooling parameter also remains unchanged, and the checkpoint contains only the registered feature layer. Reloading that state into an otherwise identical model initialized with different fixed weights changes predictions by as much as \(21.89\) in this toy. Using registered module dictionaries instead gives nonzero head and pooling updates, includes all relevant keys in the checkpoint, and reproduces predictions exactly after reload. These are actual CPU results, not predicted test outcomes.\par
The issue is Major for reproducing a trained unimodal comparator. An upstream encoder can still learn through a fixed random output projection, so the code is not necessarily unable to fit any task. Nor must a fresh reload differ if the same initialization happens to be recreated. The stronger, supported point is that the declared output heads are neither optimizer-managed nor preserved by the inspected checkpoint path, unlike the multimodal counterpart. This compromises attribution of a reproduced multimodal advantage to the training objective alone. Registration and checkpoint-roundtrip checks repair the toy mechanism; original checkpoints, benchmark effects and historical implementation remain unverified.\par
\section*{The returned continuous-behavior metric retains the trial average}
Problem 3 (Reported single-trial interpretation and metric implementation differ).  Section 4.1 says continuous decoding \(R^2\) accounts for the trial average, but gives no separate residual-metric equation. If this means removing the mean time course, the pinned evaluation does not implement it: it returns ordinary \(R^2\) after flattening all trial/time values.\par
`src/utils/eval\_utils.py`, line 762 passes raw `y` and `y\_pred` into the flattening metric. The wrapper at lines 917–943 returns this value. Its `subtract\_psth` option affects a separate plotting routine and never changes the returned scalar. The main evaluation passes `subtract="task"`, so that option’s presence does not resolve the mismatch. Training-validation selection likewise uses flattened raw values in `src/utils/utils.py`, lines 214–221.\par
For a concrete example take two observed trials \((-1,9)\) and \((1,11)\) and predict their shared average \((0,10)\) in both trials. Ordinary flattened \(R^2\) is \[1-\frac{4}{104}=\frac{25}{26}\simeq .961538.\] After subtracting the average time course, the true residuals are \((-1,-1)\) and \((1,1)\) while the predicted residuals are zero, yielding \(R^2=0\). A model can therefore score highly under the implemented metric while explaining none of this trial-to-trial variation. This does not make raw \(R^2\) an invalid metric; it makes the stronger stated interpretation unsupported by that computation. The distinction is Major for claims about residual single-trial variability. If the wording instead meant ordinary centering about one grand mean, it needs clarification rather than a change to that raw metric. The intended residual interpretation and whether the published scores used this release remain Needs Evidence.\par
A repair should specify both metrics explicitly, state how any condition-specific mean is estimated, and use the same definition across models and checkpoint selection. The example verifies the corrected residual calculation, not the paper’s scores. The test helper also places the final model call outside the batch loop in its behavioral path (lines 400–447); it retains only the last batch. Its current loader requests 10,000 examples per batch, which can conceal the issue on small sessions. This is a conditional scalability defect, not evidence that the stated IBL sessions actually crossed that limit.\par
\section*{Additional boundaries and reporting corrections}
Singleton session groups are mishandled in several routines: `argwhere(...).squeeze()` produces a scalar when a session has exactly one example. The subsequent positive-dimension guard skips session embeddings and temporal pooling weights in `encoder\_embeddings.py`, and drops such groups during validation in `trainer/base.py`, lines 333–358. An independent array-shape check confirms this case; a flat index array preserves the group. In the pooling path, skipped weights remain zero, leaving only the subsequent decoder bias for that group. This is a Minor but concrete boundary defect whose prevalence depends on batching; full-batch single-session evaluation with many trials avoids it.\par
Table 11 does not support improvement from embeddings across all five tasks. Removing temporal embeddings raises encoding from \(.20\) to \(.24\) bps, and removing both temporal and modality embeddings raises it to \(.25\); choice accuracy also rises from \(.84\) to \(.86\) in the latter row. Wheel and whisker decoding decline. The correct description is a task-dependent tradeoff, with paired session-level uncertainty needed before making significance claims. Table 2 similarly shows unchanged rounded choice accuracy when within-modality masking is removed. These are Minor reporting qualifications, not a finding that the measurements are false.\par
The roughly 400,000 neuron-hours stated on page 14 needs an exposure ledger. For session \(s\) with \(N_s\) neurons recorded for \(h_s\) hours, exposure is \(\sum_s N_s h_s\), not \((\sum_s N_s)(\sum_s h_s)\). The latter product of the headline counts, \(27{,}380\times14=383{,}320\), includes cross-session terms. A two-session check with \((N_1,N_2)=(100,200)\) and \((h_1,h_2)=(1,2)\) gives 500 neuron-hours, versus the incorrect product 900. The resemblance to the reported headline warrants reconciliation; an exact corrected dataset exposure cannot be obtained without per-session durations and neuron counts. No specific replacement total is asserted.\par
The 83\% brain-region classification result uses five-fold cross-validation over neurons from the ten already adapted animals, as Section 5.3 discloses. It measures discrimination within that pooled population, not a classifier’s generalization to a new animal excluded from all classifier folds. A leave-animal-out evaluation and controls for firing rate, session identity and embedding dimension would clarify the emergent representation claim. In particular the multimodal feature concatenates input and output weights, whereas each unimodal feature uses only one of them. Neither the confusion matrix nor predictive encoding from a single behavior identifies a causal influence or functional equivalence of brain regions. These are limits on interpretation, not demonstrated contamination of the reported classifier.\par
\section*{Evidence that survives and conclusion}
The study provides animal-level pretraining separation, several linear and transformer comparisons, extensive validation-based searches, and useful ablations. Appendix D explicitly checks a larger single-session model and reports worse performance, so the pretraining gain cannot simply be dismissed as an unexamined capacity increase. Appendix F also openly states that the POYO+ comparison uses trial-aligned windows without temporal jitter, and that NDT2 is reimplemented. The ranking is therefore for those specified adaptations and budgets. Table 10’s point estimates favor NEDS, but marginal standard errors alone do not establish paired statistical significance across ten sessions. Appendix H provides a useful separately trained monkey example, not evidence of mouse-to-monkey zero-shot transfer.\par
The supported concerns are different input conditionals for within-modality masking, unregistered unimodal output components, and a raw returned metric whose stronger residual interpretation needs clarification. Smaller batching and reporting issues also require correction. Independent toy computations establish the local mechanisms and limited repairs. They do not establish the magnitude of any historical score change or refute the value of a unified multimodal model. Original-run reconstruction, patched benchmark evaluation, raw data audits, brain-region classifier replication and author adoption remain unverified; scientific adjudication is pending.\par
1 Yizi Zhang et al., Neural Encoding and Decoding at Scale, PMLR 267 (2025), published proceedings version. https://proceedings.mlr.press/v267/zhang25bw.html.\par
\section*{Scope and limitations}
Source code supplied by the authors was not executed. The calculation is an independent finite or analytic diagnostic, not a reproduction of the original training experiments. A focused check revisited the stated Problems against the frozen original source, assumptions, decisive arguments, counterevidence and impact boundaries. No human review has occurred; scientific adjudication is pending and author adoption is unknown.\par
\end{document}
